\subsection{The first layer need not become more concentrated}
\label{sec:first-layer-scale-extension}

The finite-scale theorem also permits a conventional fixed fan-in
variance in the first layer.  This is useful because the effective
energy and its capacity bound leave the first-layer weights unrestricted.

\begin{corollary}[A fixed first-layer Gaussian scale]
\label{cor:first-layer-scale-extension}
In Theorem~\ref{thm:raw-all-scales}, replace the first-layer variance
$N\sigma^2/d=a_N^2/d$ by $v_1^2/d$, where $a_N\leq v_1\leq1$.
Leave the internal and readout priors and the likelihood unchanged.
Then the same oracle bound holds, on the same sample event, with the
same displayed constants.  In particular, one may use variance $1/d$
in the first layer for all sample sizes and widths.
\end{corollary}
\begin{proof}
The effective energy $q_N$ uses only the internal matrices and readout,
so its Gaussian exponential moment is unchanged.  The statistical event
allows arbitrary first-layer weights, and the first layer is never
modified by spectral truncation.  The prior evidence lower bound
also remains unchanged, because it uses only the centered independent
readout and bounded hidden activations.

For the teacher cylinder, first-layer noise now has normalized amplitude
$v_1\leq1$, so the bound $e_1\leq\eta(d+1)$ remains valid.
The normalized translation cost of the first-layer clone becomes
\[
 \frac{d\sigma^2}{v_1^2}\|W_*^1\|_F^2
 \leq d(d+1)B^2\frac{a_N^2}{v_1^2}
 \leq d(d+1)B^2.
\]
All other translation costs and cylinder probabilities are unchanged.
Thus the same teacher comparison applies.  The proof of
Theorem~\ref{thm:raw-all-scales} uses only these inputs.
\end{proof}
